% ============================================================
% EUSIPCO 2026 Poster
% Few-Shot Radar Person Identification via Tracking-Integrated
% InISAR and Pre-Trained Meta-Learning
% ============================================================
% Compile with: pdflatex poster.tex  (run twice)
% Requires: tikzposter class (texlive-full or scheme-full)
% ============================================================

\documentclass[a0paper,portrait]{tikzposter}

\usepackage[utf8]{inputenc}
\usepackage[T1]{fontenc}
\usepackage{amsmath,amssymb}
\usepackage{graphicx}
\usepackage{booktabs}
\usepackage{array}
\usepackage{pifont}          % for \ding{} circled numbers
\usepackage{hyperref}
\usepackage{qrcode}
\usepackage{anyfontsize}
\graphicspath{{figures/}}

% ------------------------------------------------------------
% 1. COLOR PALETTE  ("Radar Night" -- matches all diagram assets)
% ------------------------------------------------------------
\definecolor{radarNavy}{HTML}{1A2332}
\definecolor{radarPurple}{HTML}{6B4E9E}
\definecolor{radarOrange}{HTML}{E8734A}
\definecolor{radarBlue}{HTML}{3B7A9E}
\definecolor{radarGreen}{HTML}{4A9B6E}
\definecolor{radarBG}{HTML}{FAFAF8}
\definecolor{radarGrey}{HTML}{5A6472}

\colorlet{blocktitlebgcolor}{radarBlue}
\colorlet{blocktitlefgcolor}{white}
\colorlet{blockbodybgcolor}{white}
\colorlet{backgroundcolor}{radarBG}
\colorlet{framecolor}{radarNavy}

% ------------------------------------------------------------
% 2. THEME / BLOCK STYLE
% ------------------------------------------------------------
\usetheme{Default}
\usecolorstyle{Default}

% ------------------------------------------------------------
% helper command: consistent figure block with caption strip
% ------------------------------------------------------------
\newcommand{\posterfig}[3]{%
  % #1 = width, #2 = filename, #3 = caption
  \begin{center}
    \includegraphics[width=#1]{#2}\\[2mm]
    {\small\color{radarGrey}\textit{#3}}
  \end{center}
}

% ------------------------------------------------------------
% 3. TITLE
% ------------------------------------------------------------
\title{\bfseries\color{radarNavy}\parbox{0.92\linewidth}{\centering
  Few-Shot Radar Person Identification via Tracking-Integrated
  InISAR and Pre-Trained Meta-Learning}}
\author{
  Jeffee Hsiung$^{*}$ \quad Rengin Torun$^{*}$ \quad S.\ Hamed Javadi$^{\dagger}$ \quad Hichem Sahli$^{\dagger}$\\
  \normalsize KU Leuven \;|\; imec \;|\; Vrije Universiteit Brussel \\
  \footnotesize $^{*}$Equal contribution (first authors) \quad $^{\dagger}$Equal contribution (supervisors)
}
\institute{}

% ============================================================
\begin{document}
\maketitle

\begin{columns}
\column{0.5}

% ============================================================
% SECTION 1 -- MOTIVATION
% ============================================================
\block{\textcolor{white}{\large \ding{202}\ \ Why Not Just Use a Camera?}}{
\begin{minipage}{0.62\linewidth}
Conventional vision-based person identification struggles under real conditions:
\begin{itemize}
  \setlength\itemsep{2mm}
  \item \textbf{Low-light} environments degrade recognition accuracy
  \item \textbf{Adverse weather} (fog, rain) obscures visual features
  \item \textbf{Occlusions} block critical body regions
  \item \textbf{Privacy concerns} -- facial capture raises regulatory and ethical issues
\end{itemize}
\vspace{2mm}
mmWave radar avoids facial capture entirely while remaining robust to lighting, weather, and occlusion --
a privacy-preserving alternative for surveillance and human-robot interaction.
\end{minipage}
\hfill
\begin{minipage}{0.35\linewidth}
% TODO: replace with the 4-icon motivation cluster (low-light / weather / occlusion / privacy-blur)
\posterfig{\linewidth}{fig_motivation_cluster.png}{Vision fails; radar preserves privacy}
\end{minipage}
}

% ============================================================
% SECTION 2 -- WHY RADAR (sensing theory)
% ============================================================
\block{\textcolor{white}{\large \ding{203}\ \ Why mmWave Radar}}{
\begin{minipage}{0.48\linewidth}
\posterfig{\linewidth}{fig_why_radar_mimo.png}{MIMO array senses through occlusion via reflected signals}
\end{minipage}
\hfill
\begin{minipage}{0.48\linewidth}
\posterfig{\linewidth}{fig_radar_vs_camera.png}{Radar 3D point cloud (left) vs.\ camera view (right), same scene}
\end{minipage}
\vspace{3mm}

FMCW radar transmits chirps and measures round-trip time-of-flight from body-scatterers;
range and Doppler are recovered via FFT along fast-time and slow-time. This gives 3D spatial
+ velocity information \emph{without} capturing identifiable facial imagery.
}

% ============================================================
% SECTION 3 -- THE TWO PROBLEMS
% ============================================================
\block{\textcolor{white}{\large \ding{204}\ \ Two Bottlenecks in Radar-Based PID}}{
\begin{minipage}{0.48\linewidth}
\textbf{3a. Radar Imaging Quality}\\[2mm]
Conventional angle-of-arrival processing yields sparse, noisy point clouds: limited angular
resolution, motion-induced defocusing, multipath ghost artifacts (especially indoors).
InISAR has been validated on large \emph{rigid} targets (ships, aircraft) but not on small,
\emph{articulated} human targets.
\posterfig{0.95\linewidth}{fig_sparse_pointcloud_problem.png}{Sparsity, non-uniform density, noise -- missing body segments}
\posterfig{0.95\linewidth}{fig_multipath_fan.png}{Indoor multipath generates ghost clusters}
\end{minipage}
\hfill
\begin{minipage}{0.48\linewidth}
\textbf{3b. Data Scarcity}\\[2mm]
Existing learning methods need $>$100 sequences/person (hours of recording) --
impractical for adapting to new individuals. PointNet is permutation-invariant but ignores local
geometry; PCT captures local structure but targets \emph{static} objects, not temporally
evolving, noisy radar sequences.
\posterfig{0.95\linewidth}{fig_few_shot_motivation.png}{A model trained on Task 1 subjects performs near-random on unseen Task 2 subjects}
\vspace{2mm}
\textbf{Gap:} prior radar-PID + MAML work initializes meta-learning \emph{randomly}; NLP
literature shows pre-trained initialization accelerates meta-learning. We extend this to
radar point-cloud PID.
\end{minipage}
}

% ============================================================
% SECTION 4 -- OUR FRAMEWORK (3-fold contribution, system level)
% ============================================================
\block{\textcolor{white}{\large \ding{205}\ \ Our Framework: Unified Sensing + Meta-Learning}}{
\posterfig{0.92\linewidth}{01_system_architecture.png}{End-to-end system: radar signal $\to$ InISAR 3D reconstruction $\to$ spatio-temporal learning $\to$ identity}

\vspace{4mm}
\begin{minipage}{0.31\linewidth}
\centering
\colorbox{radarBlue}{\parbox{0.9\linewidth}{\centering\color{white}\textbf{(1) Radar domain}\\[1mm]
\small Tracking-integrated InISAR reconstructs high-res 3D point clouds of articulated
human motion, validated against motion-capture ground truth.}}
\end{minipage}
\hfill
\begin{minipage}{0.31\linewidth}
\centering
\colorbox{radarPurple}{\parbox{0.9\linewidth}{\centering\color{white}\textbf{(2) ML domain}\\[1mm]
\small Spatio-temporal transformer (PCTtrans) embedded in MAML; task-specific pre-trained
initialization substantially improves few-shot adaptation.}}
\end{minipage}
\hfill
\begin{minipage}{0.31\linewidth}
\centering
\colorbox{radarGreen}{\parbox{0.9\linewidth}{\centering\color{white}\textbf{(3) System integration}\\[1mm]
\small High-quality 3D reconstruction + meta-learning together enable practical few-shot
radar PID, reducing dependence on large subject-specific datasets.}}
\end{minipage}
}

% ============================================================
% COLUMN BREAK
% ============================================================
\column{0.5}

% ============================================================
% SECTION 5 -- DEEP DIVE (two sub-columns: InISAR | Learning)
% ============================================================
% \block{\textcolor{white}{\large \ding{206}a\ \ Deep Dive: InISAR Pipeline}}{
% \posterfig{\linewidth}{02_inisar_pipeline.png}{Six-stage 3D reconstruction pipeline}
% \posterfig{0.95\linewidth}{fig_motion_compensation.png}{Range recovery via FFT (fast-time/slow-time) $\to$ ISAR with Motion Compensation}
% \posterfig{0.95\linewidth}{fig_cfar_pipeline.png}{Doppler=FFT $\to$ CFAR detection $\to$ Interferometry 3D Reconstruction $\to$ Tracking}
% \posterfig{0.9\linewidth}{fig_dbscan_before_after.png}{Point-cloud preprocessing: DBSCAN $\to$ near-field filter $\to$ normalize $\to$ pad}
% }
\block{\textcolor{white}{\large \ding{206}a\ \ Deep Dive: InISAR Pipeline}}{
  \begin{minipage}{0.48\linewidth}
    \posterfig{\linewidth}{02_inisar_pipeline.png}{Six-stage 3D reconstruction pipeline}
  \end{minipage}
  \hfill
  \begin{minipage}{0.48\linewidth}
    \posterfig{\linewidth}{fig_motion_compensation.png}{Range recovery via FFT $\to$ ISAR with Motion Compensation}
  \end{minipage}

  \vspace{4mm}
  \begin{minipage}{0.48\linewidth}
    \posterfig{\linewidth}{fig_cfar_pipeline.png}{Doppler=FFT $\to$ CFAR $\to$ Interferometry 3D $\to$ Tracking}
  \end{minipage}
  \hfill
  \begin{minipage}{0.48\linewidth}
    \posterfig{\linewidth}{fig_dbscan_before_after.png}{Point-cloud preprocessing: DBSCAN $\to$ near-field filter $\to$ normalize $\to$ pad}
  \end{minipage}
}

\block{\textcolor{white}{\large \ding{206}b\ \ Deep Dive: PCTtrans + Pre-Trained MAML}}{
\posterfig{\linewidth}{04_pcttrans_highlevel.png}{High-level: spatial (per-frame) $\to$ temporal (sequence) $\to$ ID}
\posterfig{\linewidth}{fig_pcttrans_architecture.png}{PCTtrans architecture: per-frame PCT encoder $\to$ position encoding $\to$ Transformer encoder}
\vspace{4mm}
  \begin{minipage}{0.56\linewidth}
    \posterfig{\linewidth}{05_pcttrans_decomposed}{Layer decomposition: Point/Neighbor Embedding $\to$ Offset-Attention $\to$ MA-Pool $\to$ Transformer}
  \end{minipage}
  \hfill
  \begin{minipage}{0.4\linewidth}
    \posterfig{\linewidth}{fig_pretrained_maml}{Pre-trained ($\theta^{PT}_0$) vs. random ($\theta^{rand}_0$) initialization: Two-phase algorithm: supervised pre-training $\to$ MAML meta-training}
  \end{minipage}
}

% ============================================================
% SECTION 6 -- RESULTS
% ============================================================
\block{\textcolor{white}{\large \ding{207}\ \ Results}}{

\begin{minipage}{0.30\linewidth}
\centering
\colorbox{radarNavy}{\parbox{0.92\linewidth}{\centering\color{white}\vspace{1mm}
{\fontsize{40}{44}\selectfont\bfseries 98.99\%}\\[1mm]
{\small supervised accuracy, sufficient data}\vspace{1mm}}}
\end{minipage}
\hfill
\begin{minipage}{0.34\linewidth}
\centering
\colorbox{radarOrange}{\parbox{0.92\linewidth}{\centering\color{white}\vspace{1mm}
{\fontsize{34}{38}\selectfont\bfseries 32.05\%$\to$70.0\%}\\[1mm]
{\small few-shot cross-session accuracy}\vspace{1mm}}}
\end{minipage}
\hfill
\begin{minipage}{0.30\linewidth}
\centering
\colorbox{radarGreen}{\parbox{0.92\linewidth}{\centering\color{white}\vspace{1mm}
{\fontsize{40}{44}\selectfont\bfseries +15.6\,pp}\\[1mm]
{\small gain from pre-training alone}\vspace{1mm}}}
\end{minipage}

\vspace{6mm}
\begin{minipage}{0.55\linewidth}
% TODO: replace with generated bar chart (Random vs Pre-Trained init, Q1/Q2/Q3, S2)
\posterfig{\linewidth}{fig_results_barchart.png}{Random vs.\ pre-trained initialization across queries (S2)}
\end{minipage}
\hfill
\begin{minipage}{0.42\linewidth}
\small
\begin{tabular}{@{}lccc@{}}
\toprule
\textbf{Query} & \textbf{Random} & \textbf{Pre-Trained} & \textbf{Gain} \\
\midrule
Q1 (same-session)      & 76.67 & \textbf{94.44} & +17.77 \\
Q2 (cross-session)     & 54.44 & \textbf{70.00} & +15.56 \\
Q3 (cross-subject)     & 46.67 & \textbf{48.89} & +2.22 \\
\bottomrule
\end{tabular}
\vspace{3mm}

\textbf{Highlights:}
\begin{itemize}\small
  \setlength\itemsep{1mm}
  \item InISAR: 73.78\% IoU, $<$5.5\,cm height error vs.\ MoCap
  \item PCTtrans outperforms TCPCN by +2.81pp (supervised, Q1)
  \item Pre-training gain largest for cross-session (+15.56pp);
        modest for cross-subject (+2.22pp) -- unseen gaits fall
        outside pre-training distribution
\end{itemize}
\end{minipage}
}

% ============================================================
% SECTION 7 -- QR VIDEOS
% ============================================================
\block{\textcolor{white}{\large \ding{208}\ \ Watch It In Action}}{
\begin{minipage}{0.46\linewidth}
\centering
% TODO: replace URL below with your actual hosted video link
\qrcode[height=4cm]{https://example.com/inisar-reconstruction-video}\\[2mm]
\textbf{InISAR-Reconstructed 3D Point Cloud}\\
\small\color{radarGrey} Full 3D reconstruction video (ffmpeg-rendered)
\end{minipage}
\hfill
\begin{minipage}{0.46\linewidth}
\centering
% TODO: replace URL below with your actual hosted video link
\qrcode[height=4cm]{https://example.com/dbscan-before-after-video}\\[2mm]
\textbf{DBSCAN Before / After}\\
\small\color{radarGrey} Ghost-cluster removal on non-tracking-integrated point clouds
\end{minipage}
}

% ============================================================
% SECTION 8 -- INDUSTRIAL OUTLOOK / LIMITATIONS
% ============================================================
\block{\textcolor{white}{\large \ding{209}\ \ Toward Industrial Deployment: Open Challenges}}{
\small
\begin{minipage}{0.19\linewidth}
\centering
\colorbox{radarBlue!15}{\parbox{0.95\linewidth}{\vspace{1mm}
\textbf{1. Offline processing}\\[1mm]
Per-scatterer tracking is slow.\\[1mm]
\textit{Path forward:} online / cluster-center-only tracking\vspace{1mm}}}
\end{minipage}
\hfill
\begin{minipage}{0.19\linewidth}
\centering
\colorbox{radarBlue!15}{\parbox{0.95\linewidth}{\vspace{1mm}
\textbf{2. Radar-only limits}\\[1mm]
Lacks RGB-level discriminative power at scale.\\[1mm]
\textit{Path forward:} multimodal fusion where privacy allows\vspace{1mm}}}
\end{minipage}
\hfill
\begin{minipage}{0.19\linewidth}
\centering
\colorbox{radarBlue!15}{\parbox{0.95\linewidth}{\vspace{1mm}
\textbf{3. Public pretrained weights}\\[1mm]
Needs more subjects + compute.\\[1mm]
\textit{Path forward:} community dataset + HPC access\vspace{1mm}}}
\end{minipage}
\hfill
\begin{minipage}{0.19\linewidth}
\centering
\colorbox{radarBlue!15}{\parbox{0.95\linewidth}{\vspace{1mm}
\textbf{4. Tracking speed}\\[1mm]
Not yet optimized for industrial throughput.\\[1mm]
\textit{Path forward:} faster tracking algorithms\vspace{1mm}}}
\end{minipage}
\hfill
\begin{minipage}{0.19\linewidth}
\centering
\colorbox{radarBlue!15}{\parbox{0.95\linewidth}{\vspace{1mm}
\textbf{5. Fixed occlusion tolerance}\\[1mm]
Max missing frame = 1 (not adaptive).\\[1mm]
\textit{Path forward:} tunable / self-adaptive threshold\vspace{1mm}}}
\end{minipage}
}

% ============================================================
% FOOTER
% ============================================================
\block[titlewidthscale=0, bodyoffsety=0cm, roundedcorners=0]{}{
\centering\small\color{radarGrey}
jeffee.hsiung@icloud.com \quad torun10@imec.be
\qquad EUSIPCO 2026, Bruges \qquad imec / KU Leuven / VUB
}

\end{columns}
\end{document}